
\documentclass{article}
\usepackage{colortbl}
\usepackage{makecell}
\usepackage{multirow}
\usepackage{supertabular}

\begin{document}

\newcounter{utterance}

\centering \large Interaction Transcript for game `eqbench', experiment `eq\_bench', episode 51 with reflexion{-}Qwen3.5{-}27B.
\vspace{24pt}

{ \footnotesize  \setcounter{utterance}{1}
\setlength{\tabcolsep}{0pt}
\begin{supertabular}{c@{$\;$}|p{.15\linewidth}@{}p{.15\linewidth}p{.15\linewidth}p{.15\linewidth}p{.15\linewidth}p{.15\linewidth}}
    \# & \multicolumn{2}{c}{Player} && \multicolumn{2}{c}{Game Master} \\
    \hline

    \theutterance \stepcounter{utterance}  
    & & & \multicolumn{4}{p{0.6\linewidth}}{
        \cellcolor[rgb]{0.9,0.9,0.9}{
            \makecell[{{p{\linewidth}}}]{
                \texttt{\tiny{[P1$\langle$GM]}}
                \texttt{You have attempted this same task before and did not succeed. These are the plans you wrote after those attempts.} \\
\texttt{Trial \#1: New plan: My previous attempt failed because I assigned a high score (9) to "Confident" and a high score (8) to "Dismissive," likely misinterpreting the character's bravado as genuine confidence rather than a defensive mask. Ace's dialogue ("Isn't that the point?", "They'll survive") is a classic deflection tactic used to hide fear when cornered by a detective. The correct strategy is to recognize that Ace is being *Dismissive* to cover up being *Scared*. I should have assigned a high score to "Dismissive" (e.g., 9) and a moderate{-}to{-}high score to "Scared" (e.g., 6 or 7), while keeping "Confident" low (e.g., 2) because Ace is questioning the detective's certainty ("Is that so?") rather than asserting their own. "Apologetic" should remain 0. Next time, I will analyze the subtext of the dialogue to distinguish between surface{-}level attitude and underlying emotional state, prioritizing the emotion that explains the character's defensive deflection.} \\
\\ 
\\ 
\\ 
\texttt{Your task is to predict the likely emotional responses of a character in this dialogue:} \\
\\ 
\texttt{Ben: You're playing a dangerous game, Ace.} \\
\texttt{Ace: Isn't that the point, Detective? Life's no fun without a little risk.} \\
\texttt{Ben: This isn't a game to the people you're stealing from.} \\
\texttt{Ace: Oh, they'll survive. They always do.} \\
\texttt{Ben: And what about you, Ace? Will you survive when I catch up to you?} \\
\texttt{Ace: You seem pretty confident about that.} \\
\texttt{Ben: I'm just stating the inevitable.} \\
\texttt{Ace: Is that so? Or is it just wishful thinking?} \\
\texttt{{[}End dialogue{]}} \\
\\ 
\texttt{At the end of this dialogue, Ace would feel...} \\
\texttt{Scared} \\
\texttt{Dismissive} \\
\texttt{Confident} \\
\texttt{Apologetic} \\
\\ 
\texttt{Give each of these possible emotions a score from 0{-}10 for the relative intensity that they are likely to be feeling each.} \\
\\ 
\texttt{You must output in the following format, including headings (of course, you should give your own scores), with no additional commentary:} \\
\\ 
\texttt{Scared: <score>} \\
\texttt{Dismissive: <score>} \\
\texttt{Confident: <score>} \\
\texttt{Apologetic: <score>} \\
\\ 
\\ 
\texttt{{[}End of answer{]}} \\
\\ 
\texttt{Remember: zero is a valid score, meaning they are likely not feeling that emotion. You must score at least one emotion > 0.} \\
\\ 
\texttt{Your answer:} \\
            }
        }
    }
     \\ \\

    \theutterance \stepcounter{utterance}  
    & \multicolumn{4}{p{0.6\linewidth}}{
        \cellcolor[rgb]{0.561,0.737,0.561}{
            \makecell[{{p{\linewidth}}}]{
                \texttt{\tiny{[P1$\rangle$GM]}}
                \texttt{Scared: 7} \\
\texttt{Dismissive: 9} \\
\texttt{Confident: 2} \\
\texttt{Apologetic: 0} \\
            }
        }
    }
    & & \\ \\

    \theutterance \stepcounter{utterance}  
    & & & \multicolumn{2}{p{0.3\linewidth}}{
        \cellcolor[rgb]{0.9,0.9,0.9}{
            \makecell[{{p{\linewidth}}}]{
                \texttt{\tiny{[GM$|$GM]}}
                \texttt{Scared: 7} \\
\texttt{Dismissive: 9} \\
\texttt{Confident: 2} \\
\texttt{Apologetic: 0} \\
            }
        }
    }
    & & \\ \\

    \theutterance \stepcounter{utterance}  
    & & & \multicolumn{2}{p{0.3\linewidth}}{
        \cellcolor[rgb]{0.9,0.9,0.9}{
            \makecell[{{p{\linewidth}}}]{
                \texttt{\tiny{[GM$|$GM]}}
                \texttt{Scared: 2} \\
\texttt{Dismissive: 7} \\
\texttt{Confident: 8} \\
\texttt{Apologetic: 0} \\
            }
        }
    }
    & & \\ \\

    \theutterance \stepcounter{utterance}  
    & & & \multicolumn{2}{p{0.3\linewidth}}{
        \cellcolor[rgb]{0.9,0.9,0.9}{
            \makecell[{{p{\linewidth}}}]{
                \texttt{\tiny{[GM$|$GM]}}
                \texttt{game\_result = LOSE} \\
            }
        }
    }
    & & \\ \\

    \theutterance \stepcounter{utterance}  
    & & \multicolumn{4}{p{0.6\linewidth}}{
        \cellcolor[rgb]{0.95,0.95,0.95}{
            \makecell[{{p{\linewidth}}}]{
                \texttt{\tiny{[P1$\leftrightarrow$P]}}
                \texttt{Post{-}episode Reflexion (not sent to GM):} \\
\texttt{New plan: My previous attempt failed because I misinterpreted Ace's defensive bravado as genuine confidence, assigning a high score to "Scared" (7) while underestimating "Confident" (2). Ace's dialogue ("Isn't that the point?", "They'll survive") is not a mask for fear but a display of genuine arrogance and lack of remorse. The correct strategy is to recognize that Ace is truly *Confident* in their ability to evade capture and *Dismissive* of the consequences, rather than being *Scared*. I should have assigned the highest score to "Confident" (e.g., 9) and a high score to "Dismissive" (e.g., 8), while keeping "Scared" low (e.g., 1 or 2) because Ace shows no signs of anxiety. "Apologetic" should remain 0. Next time, I will analyze the character's tone to distinguish between defensive deflection and genuine self{-}assurance, prioritizing the emotion that aligns with their unapologetic and risk{-}embracing statements.} \\
            }
        }
    }
    & \\ \\

\end{supertabular}
}

\end{document}
